% Generated draft: requires mathematical and source review.
\documentclass[11pt]{article}
\usepackage[margin=25mm]{geometry}
\usepackage{fontspec}
\setmainfont{Latin Modern Roman}
\usepackage{amsmath,amssymb,mathtools,unicode-math}
\setmathfont{Latin Modern Math}
\usepackage{enumitem,needspace,xurl,hyperref}
\setlength{\parindent}{0pt}
\setlength{\parskip}{5pt}
\setlength{\emergencystretch}{3em}
\newcommand{\N}{\mathbb{N}}
\newcommand{\Z}{\mathbb{Z}}
\newcommand{\Q}{\mathbb{Q}}
\newcommand{\R}{\mathbb{R}}
\newcommand{\C}{\mathbb{C}}
\newcommand{\F}{\mathbb{F}}
\newcommand{\A}{\mathbb{A}}
\newcommand{\PP}{\mathbb{P}}
\newcommand{\E}{\mathbb{E}}
\newcommand{\eps}{\varepsilon}
\newcommand{\dd}{\,\mathrm{d}}
\begin{document}
\section*{UnsolvedMath ID: 30003946}
\textbf{Source problem code:} catalogue entry
\subsection{Definitions and mathematical statement}
Let \(X\) be a finite ideal hyperbolic simplicial complex with vertex set \(S\), and let \(\sigma,\tau\in\mathcal T(X)\). Here \(\mathcal T(X)\) denotes the space of ideal hyperbolic structures intended by the source; the supplied record does not specify its full definition or its completeness and gluing requirements. Set \(M=X\setminus S\), with domain metric \(\sigma\) and target metric \(\tau\).

Let \(\mathcal Y\) be the source's class of facewise \(W^{1,2}\) maps \(f:(M,\sigma)\to(M,\tau)\) that preserve each individual edge and face. Thus preservation means mapping an edge or face into itself, rather than merely permuting the simplices. Let
\[
\mathcal D=\{f\in\mathcal Y: f\text{ restricts to a diffeomorphism on every face}\}.
\]
The precise Sobolev convention, compatibility conditions along edges, behavior at the deleted ideal vertices, and meaning of facewise diffeomorphism at face boundaries are not specified in the supplied record and remain source-dependent.

Write \(E_{\sigma,\tau}\) for the energy functional used in the source. Its definition is absent from the supplied record. If the intended functional is the usual facewise Dirichlet energy, its conventional expression is
\[
E_{\sigma,\tau}(f)=\frac12\sum_{F}\int_{F\setminus S}|\mathrm df|_{\sigma,\tau}^{2}\,\mathrm dA_{\sigma},
\]
where the sum runs over the two-dimensional faces and the differential is measured using the domain and target metrics. This formula is an explicit conditional interpretation, not a verified reconstruction of the source's functional.

\textbf{Conjecture 2, as supplied.} For every such \(X,\sigma,\tau\), suppose that \(u\in\mathcal Y\) and \(\widetilde u\in\mathcal D\) are the respective minimizing maps, so that
\[
E_{\sigma,\tau}(u)=\inf_{f\in\mathcal Y}E_{\sigma,\tau}(f),\qquad
E_{\sigma,\tau}(\widetilde u)=\inf_{f\in\mathcal D}E_{\sigma,\tau}(f).
\]
Must
\[
u=\widetilde u
\]
hold? The source gives the equivalent formulation that the minimizing map \(u\) in \(\mathcal Y\) restricts to a diffeomorphism on every face, that is, \(u\in\mathcal D\).

\textbf{Unresolved qualification.} Membership \(u\in\mathcal D\) implies that \(u\) also minimizes over \(\mathcal D\), but equality with a separately chosen \(\widetilde u\) requires uniqueness or a specified selection of minimizers. Neither is stated in the supplied record. Equality of Sobolev maps also requires a convention concerning almost-everywhere equality or distinguished representatives. The supplied background identifies singular edges and noncompactness as unresolved difficulties; no additional regularity, compactness, or boundary hypotheses are inferred here.
\subsection{Short English statement}
For two ideal hyperbolic structures on the same finite simplicial complex, compare energy minimization among facewise Sobolev maps preserving each edge and face with minimization among those maps that are diffeomorphisms on every face. The supplied conjecture asks whether the minimizing maps coincide, and states equivalently that the Sobolev minimizer is facewise a diffeomorphism. Precise admissibility conditions, the energy definition, and uniqueness needed for this equivalence require source clarification.
\subsection{Sources}
\begin{itemize}
\item Local UnsolvedMath registry, supplied record with permanent ID \(30003946\), title \emph{Equality of Variationally Minimizing Maps}, exact source problem code \texttt{catalogue entry}. The record belongs to the Oberwolfach Reports Open Problems collection and gives proposed year \(2018\); no individual proposer is supplied.
\item Supplied report citation: \emph{New Trends in Teichm\"uller Theory and Mapping Class Groups} (2019), \emph{Oberwolfach Reports}, pp.\ 2475--2534, DOI \texttt{10.4171/owr/2018/40}. \url{https://doi.org/10.4171/owr/2018/40}. The supplied attribution identifies the conjecture as Conjecture 2.
\item Supplied literature-source entry: ``Oberwolfach Report 40, Conjecture 2,'' \url{https://doi.org/10.4171/owr/2018/40}.
\item Literature assessment dated \texttt{2026-08-21}: the supplied curation labels the problem open, reports that the source explicitly states Conjecture 2, identifies singular-edge and noncompactness difficulties, and says that no later proof was found. This is supplied dated curation, not an independently verified literature assessment.
\item The registry marks the statement \texttt{corrected\_verified} and reports that the official report was checked to restore the simplicial complex, deleted vertices, two metrics, admissible classes, and facewise-diffeomorphism formulation. These are supplied verification claims; no independent source consultation is asserted here.
\end{itemize}
\end{document}
